\section{Support of the reading at a weighted point}

% BEGIN R73P1
\paragraph{R73P1. The base character projector at the given point.}
\label{claim:R73P1}
Fix a point $q:\mathrm{WPt}(p.s,K)$, and take its own realization
$R=q.x.\mathrm{what}_2.r.z.R$, face $F=q.x.\mathrm{what}_2.r.z.F$,
coefficient field $L$ and distinguished infinite place $v_0$.
Write $G$ for its finite level group, $X$ for its good characters,
and $(P_0,\psi_0)$ for the base level of $F$.
For each $\chi\in X$, the primitive central idempotent $e_\chi$
and the level idempotent $e_{\psi_0}$ satisfy
\[
 [1](e_\chi e_{\psi_0})=\frac{\chi(1)}{|G|}\ne 0 .
\]
Indeed the two series identifications $q.\mathrm{core}.a.\mathrm{all}.\mathrm{types}.\mathrm{eSer\_def}$
identify the actual local projectors with the projectors of the given label
maps. Under the group identification in the point's $\mathrm{Pins}.\mathrm{char}$,
the good projector $\sum_{\xi\in X}e_\xi$ acts as the identity on the
absolutely irreducible $\chi$-space. In the product of the finite general
linear groups its irreducibles are tensor products. The defining good
projector selects the regular squarefree series at each factor. FLC13
and FLC17 identify the selected local constituent with the induction
of its pairwise distinct cuspidal factors. Its nondegenerate Whittaker
space has dimension one by FLC12, FLC14 and the cuspidal degree in FLC16.
The level identification in the same character pin transports the actual
$\psi_0$ to this product generic character. Taking traces of the averaging
projector and multiplying across the factors gives
\[
  \sum_{h\in P_0}\chi(h)\psi_0(h)^{-1}=|P_0|
\]
in $L$, since the equality holds after the given injective embedding
into $\mathbb C$. In the source's group-algebra conventions the
coefficient of $g$ in $e_\chi$ is $\chi(1)\chi(g^{-1})/|G|$
and $e_{\psi_0}=|P_0|^{-1}\sum_{h\in P_0}\psi_0(h)^{-1}h$.
Their product gives the asserted coefficient at the identity, with
the original character orientation.
% END R73P1

% BEGIN R73P2
\paragraph{R73P2. A nonzero row in the original base level.}
\label{claim:R73P2}
Let $a$ be a part with $E_a=R.\mathrm{EEp3}(a)\ne0$. The finite
group preserves $E_a$: its action on the row space commutes with the
type operators defining this part, by the original model and auxiliary
space identifications in $\mathrm{Pins}$ and $\mathrm{Holds}$.
On this space $\sum_{\chi\in X}e_\chi$ is the identity by the same
row model. Hence some $e_\chi E_a\ne0$, and Maschke's theorem supplies
an irreducible subrepresentation of character $\chi$ there, using the
specified coefficient embedding. On that copy $e_{\psi_0}$ has rank one
by R73P1. Any nonzero vector in its image is a row in $E_a$ transforming
under $P_0$ by $\psi_0$. The model equality identifying the level rows
with exactly those vectors puts its distinguished-place twist in
$E_a^{\mathbb C}\cap E_{\mathrm{base}}$.

For every index $i$, the original dimension equality
$\mathrm{Holds}.\mathrm{prim}.\mathrm{dimParts}(i,a)$ gives
\[
 0<\dim_{\mathbb C}(E_a^{\mathbb C}\cap E_{\mathrm{base}})
    =\dim_{\mathbb C}(W_{i,a}^{\mathbb C}\cap\operatorname{im}w_{0,i}).
\]
Here $w_{0,i}$ is the given complex base-level row map. Forgetting
the distinguished-place twist therefore gives $R.\mathrm{Wp3}(i,a)\ne0$.
This uses the dimension field already carried by $q$.
% END R73P2

% BEGIN R73P3
\paragraph{R73P3. One rational column in each nonzero part.}
\label{claim:R73P3}
Write $\iota_0$ for entrywise application of the fixed embedding at
$v_0$. For $b$ any part, put
\[
 U_b=\{y\in L^{d_W}:w_{0,i}(\iota_0 y)\in W_{i,b}^{\mathbb C}\}.
\]
The original base-level lift matrix $A$ carries $L^{d_Z}$ onto
$L^{d_W}$. Here $d_Z$ is taken at the original $i,\mu_0$.
Indeed the base rationality, $K$-form, norm and lift identities are
the consequences of the point's primitive links, character pin,
type separation and $\mathrm{Ok}$ at level zero. They give the invertible
matrix $A$ and the equality
\[
 w_{0,i}(\iota_0 Ax)=t_i(\iota_0 x)
\]
for the original $t_i=(R.\mathrm{toCx3}(F)).t_C(F.cr,\mathrm{false},i,\mu_0)$.
The rational source type spaces $Z_b$ span $L^{d_Z}$ by the
level's type identity. Its compatibility $\mathrm{H3W}$ gives
$A Z_b\subseteq U_b$. Consequently $\sum_b U_b=L^{d_W}$.

If $U_a=0$, let $M$ be the complex span of the spaces
$W_{i,b}^{\mathbb C}\cap\operatorname{im}w_{0,i}$ for $b\ne a$.
The rational decomposition just proved places $w_{0,i}(\iota_0y)$
in $M$ for every $y\in L^{d_W}$. Applying it to the coordinate basis
shows $\operatorname{im}w_{0,i}\subseteq M$. On the other hand,
$\mathrm{Holds}.\mathrm{prim}.\mathrm{sumW}(i)$ makes
$W_{i,a}^{\mathbb C}\cap\operatorname{im}w_{0,i}$ orthogonal to
each summand of $M$. Its vectors are thus orthogonal to themselves.
This contradicts R73P2. Choose $0\ne y\in U_a$ and $Ax=y$.
Then $x\ne0$ and $t_i(\iota_0 x)\in W_{i,a}^{\mathbb C}$.
The construction takes place in the original coefficient field.
% END R73P3

% BEGIN R73P4
\paragraph{R73P4. Simultaneous off-place support and preservation of the reading.}
\label{claim:R73P4}
Retain the single column $x$ of R73P3. At any infinite place $v$ the
original block hypotheses (the conjugation, irreducibility and lift
fields of $\mathrm{Holds}.\mathrm{ok}$) give a positive definite
ambient Hermitian matrix $H_v$. The base quotient projection has a
right inverse over $L$ by its stored basis identity. It remains
surjective at $v$, so the quotient metric obtained on the orthogonal
complement of its kernel is positive definite. In the original
quotient basis denote that matrix by $Q_v$. The defining frame identity
for $X_Q$ gives
\[
   \operatorname{majAt}(X_Q,v,x,x)=(\iota_vx)^*Q_v(\iota_vx)>0.
\]
Each $\iota_v$ is injective, hence $\iota_vx\ne0$. In particular the
same $x$ serves for \emph{every} $v\ne v_0$ in
$R.\mathrm{PartAt}(F,i,a,v)$: its row-membership test is always
at $v_0$, while its majorant is evaluated at $v$.
Together with R73P2 this proves all of (CSR.22) on $q$.

Apply CSR11 and CSR12 to this $R,F$ and their supplied reading.
The original zero-part update can now be made for every prescribed
re-face triple $(D_c,\mathrm{Sp},\mathrm{Zsc})$ with its two majorant
equalities and all thirteen reading identities (CSR.27).
The support step uses neither the global cuspidal multiplicity equality
nor the four joint analytic inputs carried as arguments by the upstream
\texttt{parts} wrapper. The argument starts at a supplied weighted
point and preserves its reading. Constructing the initial rational and
metric reading, and identifying its section with the requested
geometric theta, remain separate obligations.
% END R73P4

% BEGIN R73P5
\paragraph{R73P5. Removing the eventual multiplicity premise from this bundle construction.}
\label{claim:R73P5}
Consider the pinned implication
$\mathrm{GlobalTheory.Mem.R3.SelAll}.\texttt{bwit\_exists}$.
Keep its ten published inputs in their original order (Green cuspidal
characters, discrete-series cover, Harish--Chandra modules,
Casselman--Wallach, the unitary-pair Weil representation,
discrete-series existence, Harish--Chandra finiteness, cuspidal rapid
decay, global Artin reciprocity and the ray-field theorem). Keep
$h_L,N,\Gamma,p$, $p=\mathrm{prmR}(N,\mathrm{SExt.none},\Gamma)$
and the full extension premise (CSR.47). Its proof admits the following
change of hypotheses: delete
\[
 \mathrm{EvP}\bigl((\_,K)\longmapsto
       \forall q:\mathrm{WPt}(p.s,K),\;
       \mathrm{MultEq}(q.x.\mathrm{what}_2.r.z.R.\mathrm{toRealDat})\bigr).
\]
Its conclusion has the original eventual quantifiers (CSR.48): for every
such $q$, the \emph{given} $R:\mathrm{RecT}(q.O.ex)$ with
$\mathrm{RecPinE}(q,R)$, every real $c$ and every original integer
pair witnessing the two $\mathrm{Del}$ and $\mathrm{etaSz}$ equalities,
under the scale equality and $\mathrm{Val}$ guard there exists a
single $b:\mathrm{Bundle}(q.w.D(q.Z),q.O,c)$ with
$b.R=R$ and $\mathrm{Nonempty}(\mathrm{BWit}(q,b))$.

Here is the dependency reduction in the pinned proof. Start the
eventual implication with the same host/frame packet, using monotonicity
of $\mathrm{EvP}$. At $q$ take $s_P,s_S$ from $q.\mathrm{sat}$ by
transport along the given equality of columns. In the initial
selected realization, replace the call to \texttt{parts} by R73P4
on this $q$. This is the only use of the multiplicity equality in
that construction. The geometric realization transports its given
spectral-stage, extension, frame and host witnesses. Its defining
argument makes no use of this equality. The lower-map pin and the
final side realization also discard the equality on entry.

The remaining occurrence is the eigenpart construction. Its generic
character argument transports the two stored series maps to the
selected face, applies the character and group identifications there
and takes the good base idempotent, as in R73P1. For the scalar
identity, the distinguished embedding uses the realized eigenspace
and its eigenvalue. At another embedding the ray argument pulls the
part back to the old $\mathrm{mj}.eW$, uses the rational span,
family label and original $q.O.X$ scalar identification, and carries
the resulting logarithmic equality across the stored scalar-update
identity on the nonzero part. These are the actual two branches of
the reciprocity argument. Neither needs the multiplicity equality.
Consequently the same eigensystem construction takes its host,
good-character and all-embedding scalar inputs directly. The final
host datum and its pin are unaltered, so CSR21 assembles the same
$\mathrm{BWit}$ fields on this one $b$.

In the immediate linked-parameter construction
\texttt{exBL\_isLinked\_prmR}, the eventual multiplicity premise
is passed only to this bundle implication. Its other two inputs
to the $\eta$-size recursion are the movable-witness and uniform-
extension conclusions. Keeping both gives the same dependency
reduction there.

This proves a change to the hypotheses of the source constructions.
For the human proof, CSR13--22 still require their separate suppliers,
including the updated division data, geometric and spectral pins,
eigensystem outputs, all-embedding identity, host and extension
inputs. In particular we used the $\mathrm{Holds}$ and reading of a
given $q$. Other constructors that use multiplicity to produce
$\mathrm{Holds}$ on initial data retain their own hypotheses.
% END R73P5
